\section{Feature-space metrics}

\subsection{Fr\'echet distance in a learned feature space}

Let $\rho=\{\rho_i\}_{i=1}^{n}$ be a reference dataset and
$\varrho=\{\varrho_i\}_{i=1}^{m}$ a generated dataset. A feature extractor
$\psi$ maps each image to a vector representation,

\begin{equation}
    \mathbf{z}_i = \psi(\rho_i),
    \qquad
    \widetilde{\mathbf{z}}_i = \psi(\varrho_i).
\end{equation}

The two sets of features are then approximated by Gaussian distributions with
means $\boldsymbol{\mu}_{\rho}$ and $\boldsymbol{\mu}_{\varrho}$, and
covariances $\mathbf{\Sigma}_{\rho}$ and $\mathbf{\Sigma}_{\varrho}$. The
Fr\'echet distance between these two Gaussian approximations is

\begin{equation}
    d_F(\rho,\varrho)
    =
    \left\|
    \boldsymbol{\mu}_{\rho}
    -
    \boldsymbol{\mu}_{\varrho}
    \right\|_2^2
    +
    \mathrm{Tr}\!\left(
        \mathbf{\Sigma}_{\rho}
        +
        \mathbf{\Sigma}_{\varrho}
        -
        2\left(\mathbf{\Sigma}_{\rho}\mathbf{\Sigma}_{\varrho}\right)^{1/2}
    \right).
\end{equation}

Low values indicate that the two datasets occupy nearby regions of feature
space and have similar second-order statistics. This metric does not compare
images pixel by pixel; instead, it compares the distributions of datasets after
projection onto a more compact representation.

\subsection{Classical FID and PhyFID}

Two choices of feature extractor are considered in this work.

The first one is the classical FID formulation, in which $\psi$ is an
Inception-style image encoder. This provides a generic image-generation
baseline and makes it possible to check whether the generated RTI fields are
distinguishable from completely unrelated datasets or from pure noise in a
standard visual feature space.

The second one is the \emph{PhyFID} introduced for this project. Instead of
using features learned on natural images, PhyFID relies on a convolutional
autoencoder trained directly on RTI fields. The encoder part of this network
defines the mapping $\psi$, while the decoder is used only during the
reconstruction training stage. After training, each dataset is encoded into a
feature space adapted to the statistics of the considered physical domain.

In practice, the PhyFID pipeline implemented in
\texttt{lib/PhyFID/} proceeds as follows:

\begin{itemize}
    \item the training images are normalized to the interval $[0,1]$ and
    converted to tensors of shape $(N,C,H,W)$;
    \item a convolutional autoencoder is trained by minimizing a reconstruction
    mean squared error;
    \item the encoder output is used as a domain-specific feature vector;
    \item the Fr\'echet distance is finally computed between the feature
    distributions of the reference and generated datasets.
\end{itemize}

This construction is especially relevant here because the characteristic
structures of RTI fields differ strongly from the objects encountered in
standard computer vision datasets. Using a feature extractor trained on the
same family of simulations should therefore provide a more meaningful notion of
similarity for this application.


